\section{未锚定账本事件的叙事补充：northern-song}

\label{sec:anchor-batch-two-northern-song}

本组为尚无正文跳转目标的账本事件补入锚点。每一记录置于战争、行政、环境、迁徙与地方社会的关系中，并按来源限度约束叙事。

\subsection{制度、区域与材料边界}

\eventanchor{NS-1023-JIAOZI-OFFICE}{1023年·朝廷在益州设交子务，收取民间交子并发行官交子}账本条目记录宋在1023年的“朝廷在益州设交子务，收取民间交子并发行官交子”。其背景是民间交子铺户的信用危机促使地方和中央寻求可监管的票据机制，后续影响为交子纳入官府财政工具箱，纸币发行、准备金与区域流通的紧张关系长期存在。政权、军队或制度的变化须经由道路、文书、资源和地方中介才会落实，城市、乡村、边地和流动人群不可能在同一时间承受相同结果。

本条使用《宋史》卷181〈食货志下三〉；《宋会要辑稿·食货》交子条；《续资治通鉴长编》天圣元年条，并标示“官办交子事实可靠；发行额、兑现和流通范围须区分制度规定与实际运行”。这些材料能够钉住年序、命令、战役或局部地点，却不能直接推出人口、税收、身份和社会动机的总量。应区分诏令与实施、条约与边境生活、军报与伤亡统计，也不将官府的族群或行政称谓写成固定的社会本质。

从区域史角度看，该事件是资源、信息与权力重新连接的节点。不同地方会因市场、交通、宗教组织、家庭策略和环境条件而产生不同反应；材料的沉默限制我们对未被记录者所能作出的判断。

\eventanchor{NS-1023-REGNAL-YEAR}{1023年·宋以宋仁宗在位第1年作为诏令、赋役与外交文书的纪年框架}账本条目记录宋在1023年的“宋以宋仁宗在位第1年作为诏令、赋役与外交文书的纪年框架”。其背景是新岁颁历、年号和纪年是中枢与地方行政沟通的共同前提，后续影响为为同年政令、战事和地方材料提供日期锚点，不能单独推知具体政策执行。政权、军队或制度的变化须经由道路、文书、资源和地方中介才会落实，城市、乡村、边地和流动人群不可能在同一时间承受相同结果。

本条使用《宋史》卷1—23〈本纪〉；《续资治通鉴长编》本年条；《宋会要辑稿》相关门类，并标示“纪年框架可靠；该记录仅确证政权纪年连续，年内具体事项须由本年条另证”。这些材料能够钉住年序、命令、战役或局部地点，却不能直接推出人口、税收、身份和社会动机的总量。应区分诏令与实施、条约与边境生活、军报与伤亡统计，也不将官府的族群或行政称谓写成固定的社会本质。

从区域史角度看，该事件是资源、信息与权力重新连接的节点。不同地方会因市场、交通、宗教组织、家庭策略和环境条件而产生不同反应；材料的沉默限制我们对未被记录者所能作出的判断。

\eventanchor{NS-1024-REGNAL-YEAR}{1024年·宋以宋仁宗在位第2年作为诏令、赋役与外交文书的纪年框架}账本条目记录宋在1024年的“宋以宋仁宗在位第2年作为诏令、赋役与外交文书的纪年框架”。其背景是新岁颁历、年号和纪年是中枢与地方行政沟通的共同前提，后续影响为为同年政令、战事和地方材料提供日期锚点，不能单独推知具体政策执行。政权、军队或制度的变化须经由道路、文书、资源和地方中介才会落实，城市、乡村、边地和流动人群不可能在同一时间承受相同结果。

本条使用《宋史》卷1—23〈本纪〉；《续资治通鉴长编》本年条；《宋会要辑稿》相关门类，并标示“纪年框架可靠；该记录仅确证政权纪年连续，年内具体事项须由本年条另证”。这些材料能够钉住年序、命令、战役或局部地点，却不能直接推出人口、税收、身份和社会动机的总量。应区分诏令与实施、条约与边境生活、军报与伤亡统计，也不将官府的族群或行政称谓写成固定的社会本质。

从区域史角度看，该事件是资源、信息与权力重新连接的节点。不同地方会因市场、交通、宗教组织、家庭策略和环境条件而产生不同反应；材料的沉默限制我们对未被记录者所能作出的判断。

\eventanchor{NS-1025-REGNAL-YEAR}{1025年·宋以宋仁宗在位第3年作为诏令、赋役与外交文书的纪年框架}账本条目记录宋在1025年的“宋以宋仁宗在位第3年作为诏令、赋役与外交文书的纪年框架”。其背景是新岁颁历、年号和纪年是中枢与地方行政沟通的共同前提，后续影响为为同年政令、战事和地方材料提供日期锚点，不能单独推知具体政策执行。政权、军队或制度的变化须经由道路、文书、资源和地方中介才会落实，城市、乡村、边地和流动人群不可能在同一时间承受相同结果。

本条使用《宋史》卷1—23〈本纪〉；《续资治通鉴长编》本年条；《宋会要辑稿》相关门类，并标示“纪年框架可靠；该记录仅确证政权纪年连续，年内具体事项须由本年条另证”。这些材料能够钉住年序、命令、战役或局部地点，却不能直接推出人口、税收、身份和社会动机的总量。应区分诏令与实施、条约与边境生活、军报与伤亡统计，也不将官府的族群或行政称谓写成固定的社会本质。

从区域史角度看，该事件是资源、信息与权力重新连接的节点。不同地方会因市场、交通、宗教组织、家庭策略和环境条件而产生不同反应；材料的沉默限制我们对未被记录者所能作出的判断。

\eventanchor{NS-1026-REGNAL-YEAR}{1026年·宋以宋仁宗在位第4年作为诏令、赋役与外交文书的纪年框架}账本条目记录宋在1026年的“宋以宋仁宗在位第4年作为诏令、赋役与外交文书的纪年框架”。其背景是新岁颁历、年号和纪年是中枢与地方行政沟通的共同前提，后续影响为为同年政令、战事和地方材料提供日期锚点，不能单独推知具体政策执行。政权、军队或制度的变化须经由道路、文书、资源和地方中介才会落实，城市、乡村、边地和流动人群不可能在同一时间承受相同结果。

本条使用《宋史》卷1—23〈本纪〉；《续资治通鉴长编》本年条；《宋会要辑稿》相关门类，并标示“纪年框架可靠；该记录仅确证政权纪年连续，年内具体事项须由本年条另证”。这些材料能够钉住年序、命令、战役或局部地点，却不能直接推出人口、税收、身份和社会动机的总量。应区分诏令与实施、条约与边境生活、军报与伤亡统计，也不将官府的族群或行政称谓写成固定的社会本质。

从区域史角度看，该事件是资源、信息与权力重新连接的节点。不同地方会因市场、交通、宗教组织、家庭策略和环境条件而产生不同反应；材料的沉默限制我们对未被记录者所能作出的判断。

\eventanchor{NS-1027-REGNAL-YEAR}{1027年·宋以宋仁宗在位第5年作为诏令、赋役与外交文书的纪年框架}账本条目记录宋在1027年的“宋以宋仁宗在位第5年作为诏令、赋役与外交文书的纪年框架”。其背景是新岁颁历、年号和纪年是中枢与地方行政沟通的共同前提，后续影响为为同年政令、战事和地方材料提供日期锚点，不能单独推知具体政策执行。政权、军队或制度的变化须经由道路、文书、资源和地方中介才会落实，城市、乡村、边地和流动人群不可能在同一时间承受相同结果。

本条使用《宋史》卷1—23〈本纪〉；《续资治通鉴长编》本年条；《宋会要辑稿》相关门类，并标示“纪年框架可靠；该记录仅确证政权纪年连续，年内具体事项须由本年条另证”。这些材料能够钉住年序、命令、战役或局部地点，却不能直接推出人口、税收、身份和社会动机的总量。应区分诏令与实施、条约与边境生活、军报与伤亡统计，也不将官府的族群或行政称谓写成固定的社会本质。

从区域史角度看，该事件是资源、信息与权力重新连接的节点。不同地方会因市场、交通、宗教组织、家庭策略和环境条件而产生不同反应；材料的沉默限制我们对未被记录者所能作出的判断。

\eventanchor{NS-1028-REGNAL-YEAR}{1028年·宋以宋仁宗在位第6年作为诏令、赋役与外交文书的纪年框架}账本条目记录宋在1028年的“宋以宋仁宗在位第6年作为诏令、赋役与外交文书的纪年框架”。其背景是新岁颁历、年号和纪年是中枢与地方行政沟通的共同前提，后续影响为为同年政令、战事和地方材料提供日期锚点，不能单独推知具体政策执行。政权、军队或制度的变化须经由道路、文书、资源和地方中介才会落实，城市、乡村、边地和流动人群不可能在同一时间承受相同结果。

本条使用《宋史》卷1—23〈本纪〉；《续资治通鉴长编》本年条；《宋会要辑稿》相关门类，并标示“纪年框架可靠；该记录仅确证政权纪年连续，年内具体事项须由本年条另证”。这些材料能够钉住年序、命令、战役或局部地点，却不能直接推出人口、税收、身份和社会动机的总量。应区分诏令与实施、条约与边境生活、军报与伤亡统计，也不将官府的族群或行政称谓写成固定的社会本质。

从区域史角度看，该事件是资源、信息与权力重新连接的节点。不同地方会因市场、交通、宗教组织、家庭策略和环境条件而产生不同反应；材料的沉默限制我们对未被记录者所能作出的判断。

\eventanchor{NS-1029-REGNAL-YEAR}{1029年·宋以宋仁宗在位第7年作为诏令、赋役与外交文书的纪年框架}账本条目记录宋在1029年的“宋以宋仁宗在位第7年作为诏令、赋役与外交文书的纪年框架”。其背景是新岁颁历、年号和纪年是中枢与地方行政沟通的共同前提，后续影响为为同年政令、战事和地方材料提供日期锚点，不能单独推知具体政策执行。政权、军队或制度的变化须经由道路、文书、资源和地方中介才会落实，城市、乡村、边地和流动人群不可能在同一时间承受相同结果。

本条使用《宋史》卷1—23〈本纪〉；《续资治通鉴长编》本年条；《宋会要辑稿》相关门类，并标示“纪年框架可靠；该记录仅确证政权纪年连续，年内具体事项须由本年条另证”。这些材料能够钉住年序、命令、战役或局部地点，却不能直接推出人口、税收、身份和社会动机的总量。应区分诏令与实施、条约与边境生活、军报与伤亡统计，也不将官府的族群或行政称谓写成固定的社会本质。

从区域史角度看，该事件是资源、信息与权力重新连接的节点。不同地方会因市场、交通、宗教组织、家庭策略和环境条件而产生不同反应；材料的沉默限制我们对未被记录者所能作出的判断。

\eventanchor{NS-1030-REGNAL-YEAR}{1030年·宋以宋仁宗在位第8年作为诏令、赋役与外交文书的纪年框架}账本条目记录宋在1030年的“宋以宋仁宗在位第8年作为诏令、赋役与外交文书的纪年框架”。其背景是新岁颁历、年号和纪年是中枢与地方行政沟通的共同前提，后续影响为为同年政令、战事和地方材料提供日期锚点，不能单独推知具体政策执行。政权、军队或制度的变化须经由道路、文书、资源和地方中介才会落实，城市、乡村、边地和流动人群不可能在同一时间承受相同结果。

本条使用《宋史》卷1—23〈本纪〉；《续资治通鉴长编》本年条；《宋会要辑稿》相关门类，并标示“纪年框架可靠；该记录仅确证政权纪年连续，年内具体事项须由本年条另证”。这些材料能够钉住年序、命令、战役或局部地点，却不能直接推出人口、税收、身份和社会动机的总量。应区分诏令与实施、条约与边境生活、军报与伤亡统计，也不将官府的族群或行政称谓写成固定的社会本质。

从区域史角度看，该事件是资源、信息与权力重新连接的节点。不同地方会因市场、交通、宗教组织、家庭策略和环境条件而产生不同反应；材料的沉默限制我们对未被记录者所能作出的判断。

\eventanchor{NS-1031-REGNAL-YEAR}{1031年·宋以宋仁宗在位第9年作为诏令、赋役与外交文书的纪年框架}账本条目记录宋在1031年的“宋以宋仁宗在位第9年作为诏令、赋役与外交文书的纪年框架”。其背景是新岁颁历、年号和纪年是中枢与地方行政沟通的共同前提，后续影响为为同年政令、战事和地方材料提供日期锚点，不能单独推知具体政策执行。政权、军队或制度的变化须经由道路、文书、资源和地方中介才会落实，城市、乡村、边地和流动人群不可能在同一时间承受相同结果。

本条使用《宋史》卷1—23〈本纪〉；《续资治通鉴长编》本年条；《宋会要辑稿》相关门类，并标示“纪年框架可靠；该记录仅确证政权纪年连续，年内具体事项须由本年条另证”。这些材料能够钉住年序、命令、战役或局部地点，却不能直接推出人口、税收、身份和社会动机的总量。应区分诏令与实施、条约与边境生活、军报与伤亡统计，也不将官府的族群或行政称谓写成固定的社会本质。

从区域史角度看，该事件是资源、信息与权力重新连接的节点。不同地方会因市场、交通、宗教组织、家庭策略和环境条件而产生不同反应；材料的沉默限制我们对未被记录者所能作出的判断。

\eventanchor{NS-1032-REGNAL-YEAR}{1032年·宋以宋仁宗在位第10年作为诏令、赋役与外交文书的纪年框架}账本条目记录宋在1032年的“宋以宋仁宗在位第10年作为诏令、赋役与外交文书的纪年框架”。其背景是新岁颁历、年号和纪年是中枢与地方行政沟通的共同前提，后续影响为为同年政令、战事和地方材料提供日期锚点，不能单独推知具体政策执行。政权、军队或制度的变化须经由道路、文书、资源和地方中介才会落实，城市、乡村、边地和流动人群不可能在同一时间承受相同结果。

本条使用《宋史》卷1—23〈本纪〉；《续资治通鉴长编》本年条；《宋会要辑稿》相关门类，并标示“纪年框架可靠；该记录仅确证政权纪年连续，年内具体事项须由本年条另证”。这些材料能够钉住年序、命令、战役或局部地点，却不能直接推出人口、税收、身份和社会动机的总量。应区分诏令与实施、条约与边境生活、军报与伤亡统计，也不将官府的族群或行政称谓写成固定的社会本质。

从区域史角度看，该事件是资源、信息与权力重新连接的节点。不同地方会因市场、交通、宗教组织、家庭策略和环境条件而产生不同反应；材料的沉默限制我们对未被记录者所能作出的判断。

\eventanchor{NS-1033-REGNAL-YEAR}{1033年·宋以宋仁宗在位第11年作为诏令、赋役与外交文书的纪年框架}账本条目记录宋在1033年的“宋以宋仁宗在位第11年作为诏令、赋役与外交文书的纪年框架”。其背景是新岁颁历、年号和纪年是中枢与地方行政沟通的共同前提，后续影响为为同年政令、战事和地方材料提供日期锚点，不能单独推知具体政策执行。政权、军队或制度的变化须经由道路、文书、资源和地方中介才会落实，城市、乡村、边地和流动人群不可能在同一时间承受相同结果。

本条使用《宋史》卷1—23〈本纪〉；《续资治通鉴长编》本年条；《宋会要辑稿》相关门类，并标示“纪年框架可靠；该记录仅确证政权纪年连续，年内具体事项须由本年条另证”。这些材料能够钉住年序、命令、战役或局部地点，却不能直接推出人口、税收、身份和社会动机的总量。应区分诏令与实施、条约与边境生活、军报与伤亡统计，也不将官府的族群或行政称谓写成固定的社会本质。

从区域史角度看，该事件是资源、信息与权力重新连接的节点。不同地方会因市场、交通、宗教组织、家庭策略和环境条件而产生不同反应；材料的沉默限制我们对未被记录者所能作出的判断。

\eventanchor{NS-1033-RENZONG-PERSONAL-RULE}{1033年·刘太后去世后仁宗开始亲政}账本条目记录宋在1033年的“刘太后去世后仁宗开始亲政”。其背景是长期摄政结束，需要重排太后旧臣、台谏与宰执之间的关系，后续影响为仁宗朝的政治议论空间扩大，并在西夏危机下接受庆历改革议程。政权、军队或制度的变化须经由道路、文书、资源和地方中介才会落实，城市、乡村、边地和流动人群不可能在同一时间承受相同结果。

本条使用《宋史》卷8〈仁宗本纪一〉；《续资治通鉴长编》明道二年条；《宋会要辑稿·职官》宰辅条，并标示“政权转换可靠；所谓完全由皇帝立即掌政忽略辅臣和内廷网络”。这些材料能够钉住年序、命令、战役或局部地点，却不能直接推出人口、税收、身份和社会动机的总量。应区分诏令与实施、条约与边境生活、军报与伤亡统计，也不将官府的族群或行政称谓写成固定的社会本质。

从区域史角度看，该事件是资源、信息与权力重新连接的节点。不同地方会因市场、交通、宗教组织、家庭策略和环境条件而产生不同反应；材料的沉默限制我们对未被记录者所能作出的判断。

\eventanchor{NS-1034-REGNAL-YEAR}{1034年·宋以宋仁宗在位第12年作为诏令、赋役与外交文书的纪年框架}账本条目记录宋在1034年的“宋以宋仁宗在位第12年作为诏令、赋役与外交文书的纪年框架”。其背景是新岁颁历、年号和纪年是中枢与地方行政沟通的共同前提，后续影响为为同年政令、战事和地方材料提供日期锚点，不能单独推知具体政策执行。政权、军队或制度的变化须经由道路、文书、资源和地方中介才会落实，城市、乡村、边地和流动人群不可能在同一时间承受相同结果。

本条使用《宋史》卷1—23〈本纪〉；《续资治通鉴长编》本年条；《宋会要辑稿》相关门类，并标示“纪年框架可靠；该记录仅确证政权纪年连续，年内具体事项须由本年条另证”。这些材料能够钉住年序、命令、战役或局部地点，却不能直接推出人口、税收、身份和社会动机的总量。应区分诏令与实施、条约与边境生活、军报与伤亡统计，也不将官府的族群或行政称谓写成固定的社会本质。

从区域史角度看，该事件是资源、信息与权力重新连接的节点。不同地方会因市场、交通、宗教组织、家庭策略和环境条件而产生不同反应；材料的沉默限制我们对未被记录者所能作出的判断。

\eventanchor{NS-1034-YUANHAO-BREAK}{1034年·李元昊改用新国号和年号，逐步脱离宋册封秩序}账本条目记录宋与党项在1034年的“李元昊改用新国号和年号，逐步脱离宋册封秩序”。其背景是党项首领整合河套、河西资源并重塑王权象征，与宋的册封关系日益紧张，后续影响为宋夏冲突从边境摩擦升级为政权间战争，西夏自身制度应另取西夏文材料核验。政权、军队或制度的变化须经由道路、文书、资源和地方中介才会落实，城市、乡村、边地和流动人群不可能在同一时间承受相同结果。

本条使用《宋史》卷485〈夏国传上〉；《续资治通鉴长编》景祐元年条；《宋会要辑稿·蕃夷》夏国条，并标示“改号事实可靠；宋方贬称与党项内部政治动机不能混为一谈”。这些材料能够钉住年序、命令、战役或局部地点，却不能直接推出人口、税收、身份和社会动机的总量。应区分诏令与实施、条约与边境生活、军报与伤亡统计，也不将官府的族群或行政称谓写成固定的社会本质。

从区域史角度看，该事件是资源、信息与权力重新连接的节点。不同地方会因市场、交通、宗教组织、家庭策略和环境条件而产生不同反应；材料的沉默限制我们对未被记录者所能作出的判断。

\eventanchor{NS-1035-REGNAL-YEAR}{1035年·宋以宋仁宗在位第13年作为诏令、赋役与外交文书的纪年框架}账本条目记录宋在1035年的“宋以宋仁宗在位第13年作为诏令、赋役与外交文书的纪年框架”。其背景是新岁颁历、年号和纪年是中枢与地方行政沟通的共同前提，后续影响为为同年政令、战事和地方材料提供日期锚点，不能单独推知具体政策执行。政权、军队或制度的变化须经由道路、文书、资源和地方中介才会落实，城市、乡村、边地和流动人群不可能在同一时间承受相同结果。

本条使用《宋史》卷1—23〈本纪〉；《续资治通鉴长编》本年条；《宋会要辑稿》相关门类，并标示“纪年框架可靠；该记录仅确证政权纪年连续，年内具体事项须由本年条另证”。这些材料能够钉住年序、命令、战役或局部地点，却不能直接推出人口、税收、身份和社会动机的总量。应区分诏令与实施、条约与边境生活、军报与伤亡统计，也不将官府的族群或行政称谓写成固定的社会本质。

从区域史角度看，该事件是资源、信息与权力重新连接的节点。不同地方会因市场、交通、宗教组织、家庭策略和环境条件而产生不同反应；材料的沉默限制我们对未被记录者所能作出的判断。

\eventanchor{NS-1036-REGNAL-YEAR}{1036年·宋以宋仁宗在位第14年作为诏令、赋役与外交文书的纪年框架}账本条目记录宋在1036年的“宋以宋仁宗在位第14年作为诏令、赋役与外交文书的纪年框架”。其背景是新岁颁历、年号和纪年是中枢与地方行政沟通的共同前提，后续影响为为同年政令、战事和地方材料提供日期锚点，不能单独推知具体政策执行。政权、军队或制度的变化须经由道路、文书、资源和地方中介才会落实，城市、乡村、边地和流动人群不可能在同一时间承受相同结果。

本条使用《宋史》卷1—23〈本纪〉；《续资治通鉴长编》本年条；《宋会要辑稿》相关门类，并标示“纪年框架可靠；该记录仅确证政权纪年连续，年内具体事项须由本年条另证”。这些材料能够钉住年序、命令、战役或局部地点，却不能直接推出人口、税收、身份和社会动机的总量。应区分诏令与实施、条约与边境生活、军报与伤亡统计，也不将官府的族群或行政称谓写成固定的社会本质。

从区域史角度看，该事件是资源、信息与权力重新连接的节点。不同地方会因市场、交通、宗教组织、家庭策略和环境条件而产生不同反应；材料的沉默限制我们对未被记录者所能作出的判断。

\eventanchor{NS-1037-REGNAL-YEAR}{1037年·宋以宋仁宗在位第15年作为诏令、赋役与外交文书的纪年框架}账本条目记录宋在1037年的“宋以宋仁宗在位第15年作为诏令、赋役与外交文书的纪年框架”。其背景是新岁颁历、年号和纪年是中枢与地方行政沟通的共同前提，后续影响为为同年政令、战事和地方材料提供日期锚点，不能单独推知具体政策执行。政权、军队或制度的变化须经由道路、文书、资源和地方中介才会落实，城市、乡村、边地和流动人群不可能在同一时间承受相同结果。

本条使用《宋史》卷1—23〈本纪〉；《续资治通鉴长编》本年条；《宋会要辑稿》相关门类，并标示“纪年框架可靠；该记录仅确证政权纪年连续，年内具体事项须由本年条另证”。这些材料能够钉住年序、命令、战役或局部地点，却不能直接推出人口、税收、身份和社会动机的总量。应区分诏令与实施、条约与边境生活、军报与伤亡统计，也不将官府的族群或行政称谓写成固定的社会本质。

从区域史角度看，该事件是资源、信息与权力重新连接的节点。不同地方会因市场、交通、宗教组织、家庭策略和环境条件而产生不同反应；材料的沉默限制我们对未被记录者所能作出的判断。

\eventanchor{NS-1038-REGNAL-YEAR}{1038年·宋以宋仁宗在位第16年作为诏令、赋役与外交文书的纪年框架}账本条目记录宋在1038年的“宋以宋仁宗在位第16年作为诏令、赋役与外交文书的纪年框架”。其背景是新岁颁历、年号和纪年是中枢与地方行政沟通的共同前提，后续影响为为同年政令、战事和地方材料提供日期锚点，不能单独推知具体政策执行。政权、军队或制度的变化须经由道路、文书、资源和地方中介才会落实，城市、乡村、边地和流动人群不可能在同一时间承受相同结果。

本条使用《宋史》卷1—23〈本纪〉；《续资治通鉴长编》本年条；《宋会要辑稿》相关门类，并标示“纪年框架可靠；该记录仅确证政权纪年连续，年内具体事项须由本年条另证”。这些材料能够钉住年序、命令、战役或局部地点，却不能直接推出人口、税收、身份和社会动机的总量。应区分诏令与实施、条约与边境生活、军报与伤亡统计，也不将官府的族群或行政称谓写成固定的社会本质。

从区域史角度看，该事件是资源、信息与权力重新连接的节点。不同地方会因市场、交通、宗教组织、家庭策略和环境条件而产生不同反应；材料的沉默限制我们对未被记录者所能作出的判断。

\eventanchor{NS-1038-WESTERN-XIA-FOUNDATION}{1038年·李元昊称帝，西夏政权正式建立}账本条目记录西夏与宋在1038年的“李元昊称帝，西夏政权正式建立”。其背景是李元昊已完成权力集中，称帝以改变对宋辽的外交位阶，后续影响为宋夏战争加剧，西夏国家史不能仅由宋朝战报和朝贡术语界定。政权、军队或制度的变化须经由道路、文书、资源和地方中介才会落实，城市、乡村、边地和流动人群不可能在同一时间承受相同结果。

本条使用《宋史》卷485〈夏国传上〉；《续资治通鉴长编》宝元元年条；《宋会要辑稿·蕃夷》夏国条，并标示“称帝可靠；建国礼制和内部接受程度以宋方记录为主，须保留外部观察限制”。这些材料能够钉住年序、命令、战役或局部地点，却不能直接推出人口、税收、身份和社会动机的总量。应区分诏令与实施、条约与边境生活、军报与伤亡统计，也不将官府的族群或行政称谓写成固定的社会本质。

从区域史角度看，该事件是资源、信息与权力重新连接的节点。不同地方会因市场、交通、宗教组织、家庭策略和环境条件而产生不同反应；材料的沉默限制我们对未被记录者所能作出的判断。

\eventanchor{NS-1039-REGNAL-YEAR}{1039年·宋以宋仁宗在位第17年作为诏令、赋役与外交文书的纪年框架}账本条目记录宋在1039年的“宋以宋仁宗在位第17年作为诏令、赋役与外交文书的纪年框架”。其背景是新岁颁历、年号和纪年是中枢与地方行政沟通的共同前提，后续影响为为同年政令、战事和地方材料提供日期锚点，不能单独推知具体政策执行。政权、军队或制度的变化须经由道路、文书、资源和地方中介才会落实，城市、乡村、边地和流动人群不可能在同一时间承受相同结果。

本条使用《宋史》卷1—23〈本纪〉；《续资治通鉴长编》本年条；《宋会要辑稿》相关门类，并标示“纪年框架可靠；该记录仅确证政权纪年连续，年内具体事项须由本年条另证”。这些材料能够钉住年序、命令、战役或局部地点，却不能直接推出人口、税收、身份和社会动机的总量。应区分诏令与实施、条约与边境生活、军报与伤亡统计，也不将官府的族群或行政称谓写成固定的社会本质。

从区域史角度看，该事件是资源、信息与权力重新连接的节点。不同地方会因市场、交通、宗教组织、家庭策略和环境条件而产生不同反应；材料的沉默限制我们对未被记录者所能作出的判断。

\eventanchor{NS-1040-REGNAL-YEAR}{1040年·宋以宋仁宗在位第18年作为诏令、赋役与外交文书的纪年框架}账本条目记录宋在1040年的“宋以宋仁宗在位第18年作为诏令、赋役与外交文书的纪年框架”。其背景是新岁颁历、年号和纪年是中枢与地方行政沟通的共同前提，后续影响为为同年政令、战事和地方材料提供日期锚点，不能单独推知具体政策执行。政权、军队或制度的变化须经由道路、文书、资源和地方中介才会落实，城市、乡村、边地和流动人群不可能在同一时间承受相同结果。

本条使用《宋史》卷1—23〈本纪〉；《续资治通鉴长编》本年条；《宋会要辑稿》相关门类，并标示“纪年框架可靠；该记录仅确证政权纪年连续，年内具体事项须由本年条另证”。这些材料能够钉住年序、命令、战役或局部地点，却不能直接推出人口、税收、身份和社会动机的总量。应区分诏令与实施、条约与边境生活、军报与伤亡统计，也不将官府的族群或行政称谓写成固定的社会本质。

从区域史角度看，该事件是资源、信息与权力重新连接的节点。不同地方会因市场、交通、宗教组织、家庭策略和环境条件而产生不同反应；材料的沉默限制我们对未被记录者所能作出的判断。

\eventanchor{NS-1040-SANCHUANKOU}{1040年·三川口之战中宋军失利}账本条目记录宋与西夏在1040年的“三川口之战中宋军失利”。其背景是宋军试图在陕北前线压制西夏，行军和补给受黄土高原地形制约，后续影响为宋廷调整西北防务、募兵和边费安排，边境战争拖入持久消耗。政权、军队或制度的变化须经由道路、文书、资源和地方中介才会落实，城市、乡村、边地和流动人群不可能在同一时间承受相同结果。

本条使用《宋史》卷9〈仁宗本纪二〉；《宋史》卷318〈范雍传〉；《续资治通鉴长编》康定元年条，并标示“宋军败绩可靠；兵力、伤亡和将领责任应并列宋方与西夏研究”。这些材料能够钉住年序、命令、战役或局部地点，却不能直接推出人口、税收、身份和社会动机的总量。应区分诏令与实施、条约与边境生活、军报与伤亡统计，也不将官府的族群或行政称谓写成固定的社会本质。

从区域史角度看，该事件是资源、信息与权力重新连接的节点。不同地方会因市场、交通、宗教组织、家庭策略和环境条件而产生不同反应；材料的沉默限制我们对未被记录者所能作出的判断。

\eventanchor{NS-1041-HAOSHUI-CHUAN}{1041年·好水川之战中宋军再遭挫败}账本条目记录宋与西夏在1041年的“好水川之战中宋军再遭挫败”。其背景是宋军在西夏机动作战与山地补给压力下寻求决战，后续影响为连续失利削弱速胜预期，促使宋廷转向堡寨、和议和边防财政的组合。政权、军队或制度的变化须经由道路、文书、资源和地方中介才会落实，城市、乡村、边地和流动人群不可能在同一时间承受相同结果。

本条使用《宋史》卷9〈仁宗本纪二〉；《宋史》卷314〈任福传〉；《续资治通鉴长编》庆历元年条，并标示“战败可靠；诱敌叙事和死伤数主要来自宋方编年，须谨慎”。这些材料能够钉住年序、命令、战役或局部地点，却不能直接推出人口、税收、身份和社会动机的总量。应区分诏令与实施、条约与边境生活、军报与伤亡统计，也不将官府的族群或行政称谓写成固定的社会本质。

从区域史角度看，该事件是资源、信息与权力重新连接的节点。不同地方会因市场、交通、宗教组织、家庭策略和环境条件而产生不同反应；材料的沉默限制我们对未被记录者所能作出的判断。

\eventanchor{NS-1041-REGNAL-YEAR}{1041年·宋以宋仁宗在位第19年作为诏令、赋役与外交文书的纪年框架}账本条目记录宋在1041年的“宋以宋仁宗在位第19年作为诏令、赋役与外交文书的纪年框架”。其背景是新岁颁历、年号和纪年是中枢与地方行政沟通的共同前提，后续影响为为同年政令、战事和地方材料提供日期锚点，不能单独推知具体政策执行。政权、军队或制度的变化须经由道路、文书、资源和地方中介才会落实，城市、乡村、边地和流动人群不可能在同一时间承受相同结果。

本条使用《宋史》卷1—23〈本纪〉；《续资治通鉴长编》本年条；《宋会要辑稿》相关门类，并标示“纪年框架可靠；该记录仅确证政权纪年连续，年内具体事项须由本年条另证”。这些材料能够钉住年序、命令、战役或局部地点，却不能直接推出人口、税收、身份和社会动机的总量。应区分诏令与实施、条约与边境生活、军报与伤亡统计，也不将官府的族群或行政称谓写成固定的社会本质。

从区域史角度看，该事件是资源、信息与权力重新连接的节点。不同地方会因市场、交通、宗教组织、家庭策略和环境条件而产生不同反应；材料的沉默限制我们对未被记录者所能作出的判断。

\subsection{制度、区域与材料边界}

\eventanchor{NS-1042-DINGCHUAN}{1042年·定川寨之战后宋廷西北防务进一步承压}账本条目记录宋与西夏在1042年的“定川寨之战后宋廷西北防务进一步承压”。其背景是宋夏边界缺乏稳定缓冲，双方围绕寨堡、道路和水源反复争夺，后续影响为西夏战争成本推动宋廷同时讨论改革吏治、财政调度和外交妥协。政权、军队或制度的变化须经由道路、文书、资源和地方中介才会落实，城市、乡村、边地和流动人群不可能在同一时间承受相同结果。

本条使用《宋史》卷10〈仁宗本纪三〉；《宋史》卷485〈夏国传上〉；《续资治通鉴长编》庆历二年条，并标示“战事可靠；地望、军数与战果须与考古地理研究互校”。这些材料能够钉住年序、命令、战役或局部地点，却不能直接推出人口、税收、身份和社会动机的总量。应区分诏令与实施、条约与边境生活、军报与伤亡统计，也不将官府的族群或行政称谓写成固定的社会本质。

从区域史角度看，该事件是资源、信息与权力重新连接的节点。不同地方会因市场、交通、宗教组织、家庭策略和环境条件而产生不同反应；材料的沉默限制我们对未被记录者所能作出的判断。

\eventanchor{NS-1042-REGNAL-YEAR}{1042年·宋以宋仁宗在位第20年作为诏令、赋役与外交文书的纪年框架}账本条目记录宋在1042年的“宋以宋仁宗在位第20年作为诏令、赋役与外交文书的纪年框架”。其背景是新岁颁历、年号和纪年是中枢与地方行政沟通的共同前提，后续影响为为同年政令、战事和地方材料提供日期锚点，不能单独推知具体政策执行。政权、军队或制度的变化须经由道路、文书、资源和地方中介才会落实，城市、乡村、边地和流动人群不可能在同一时间承受相同结果。

本条使用《宋史》卷1—23〈本纪〉；《续资治通鉴长编》本年条；《宋会要辑稿》相关门类，并标示“纪年框架可靠；该记录仅确证政权纪年连续，年内具体事项须由本年条另证”。这些材料能够钉住年序、命令、战役或局部地点，却不能直接推出人口、税收、身份和社会动机的总量。应区分诏令与实施、条约与边境生活、军报与伤亡统计，也不将官府的族群或行政称谓写成固定的社会本质。

从区域史角度看，该事件是资源、信息与权力重新连接的节点。不同地方会因市场、交通、宗教组织、家庭策略和环境条件而产生不同反应；材料的沉默限制我们对未被记录者所能作出的判断。

\eventanchor{NS-1043-QINGLI-REFORMS}{1043年·范仲淹、富弼等提出并推行庆历新政的一组吏治与学校措施}账本条目记录宋在1043年的“范仲淹、富弼等提出并推行庆历新政的一组吏治与学校措施”。其背景是西夏战事、冗官财政和官僚选任问题促使中枢寻求整饬方案，后续影响为改革激化官僚联盟冲突，成为后续熙宁变法争论的重要参照而非直接前身。政权、军队或制度的变化须经由道路、文书、资源和地方中介才会落实，城市、乡村、边地和流动人群不可能在同一时间承受相同结果。

本条使用《宋史》卷10〈仁宗本纪三〉；《宋会要辑稿·职官》庆历条；《续资治通鉴长编》庆历三年条，并标示“政策提出可靠；各项推行程度和地方效果不能由诏令直接推定”。这些材料能够钉住年序、命令、战役或局部地点，却不能直接推出人口、税收、身份和社会动机的总量。应区分诏令与实施、条约与边境生活、军报与伤亡统计，也不将官府的族群或行政称谓写成固定的社会本质。

从区域史角度看，该事件是资源、信息与权力重新连接的节点。不同地方会因市场、交通、宗教组织、家庭策略和环境条件而产生不同反应；材料的沉默限制我们对未被记录者所能作出的判断。

\eventanchor{NS-1043-REGNAL-YEAR}{1043年·宋以宋仁宗在位第21年作为诏令、赋役与外交文书的纪年框架}账本条目记录宋在1043年的“宋以宋仁宗在位第21年作为诏令、赋役与外交文书的纪年框架”。其背景是新岁颁历、年号和纪年是中枢与地方行政沟通的共同前提，后续影响为为同年政令、战事和地方材料提供日期锚点，不能单独推知具体政策执行。政权、军队或制度的变化须经由道路、文书、资源和地方中介才会落实，城市、乡村、边地和流动人群不可能在同一时间承受相同结果。

本条使用《宋史》卷1—23〈本纪〉；《续资治通鉴长编》本年条；《宋会要辑稿》相关门类，并标示“纪年框架可靠；该记录仅确证政权纪年连续，年内具体事项须由本年条另证”。这些材料能够钉住年序、命令、战役或局部地点，却不能直接推出人口、税收、身份和社会动机的总量。应区分诏令与实施、条约与边境生活、军报与伤亡统计，也不将官府的族群或行政称谓写成固定的社会本质。

从区域史角度看，该事件是资源、信息与权力重新连接的节点。不同地方会因市场、交通、宗教组织、家庭策略和环境条件而产生不同反应；材料的沉默限制我们对未被记录者所能作出的判断。

\eventanchor{NS-1044-REGNAL-YEAR}{1044年·宋以宋仁宗在位第22年作为诏令、赋役与外交文书的纪年框架}账本条目记录宋在1044年的“宋以宋仁宗在位第22年作为诏令、赋役与外交文书的纪年框架”。其背景是新岁颁历、年号和纪年是中枢与地方行政沟通的共同前提，后续影响为为同年政令、战事和地方材料提供日期锚点，不能单独推知具体政策执行。政权、军队或制度的变化须经由道路、文书、资源和地方中介才会落实，城市、乡村、边地和流动人群不可能在同一时间承受相同结果。

本条使用《宋史》卷1—23〈本纪〉；《续资治通鉴长编》本年条；《宋会要辑稿》相关门类，并标示“纪年框架可靠；该记录仅确证政权纪年连续，年内具体事项须由本年条另证”。这些材料能够钉住年序、命令、战役或局部地点，却不能直接推出人口、税收、身份和社会动机的总量。应区分诏令与实施、条约与边境生活、军报与伤亡统计，也不将官府的族群或行政称谓写成固定的社会本质。

从区域史角度看，该事件是资源、信息与权力重新连接的节点。不同地方会因市场、交通、宗教组织、家庭策略和环境条件而产生不同反应；材料的沉默限制我们对未被记录者所能作出的判断。

\eventanchor{NS-1044-SONG-XIA-PEACE}{1044年·宋夏达成庆历和议，宋以赐予和互市安排换取边境停战}账本条目记录宋与西夏在1044年的“宋夏达成庆历和议，宋以赐予和互市安排换取边境停战”。其背景是连续边战未能使任何一方取得决定性控制，双方皆承受补给压力，后续影响为西北获得相对稳定，宋廷可重新配置边费；边境贸易和寨堡竞争仍未消失。政权、军队或制度的变化须经由道路、文书、资源和地方中介才会落实，城市、乡村、边地和流动人群不可能在同一时间承受相同结果。

本条使用《宋史》卷10〈仁宗本纪三〉；《宋史》卷485〈夏国传上〉；《宋会要辑稿·蕃夷》夏国和议条，并标示“和议可靠；赐予并非单向臣属证明，条款和执行须与双方称谓分开”。这些材料能够钉住年序、命令、战役或局部地点，却不能直接推出人口、税收、身份和社会动机的总量。应区分诏令与实施、条约与边境生活、军报与伤亡统计，也不将官府的族群或行政称谓写成固定的社会本质。

从区域史角度看，该事件是资源、信息与权力重新连接的节点。不同地方会因市场、交通、宗教组织、家庭策略和环境条件而产生不同反应；材料的沉默限制我们对未被记录者所能作出的判断。

\eventanchor{NS-1045-QINGLI-REFORM-END}{1045年·庆历新政的主要倡议者相继离开中枢，改革推进中止}账本条目记录宋在1045年的“庆历新政的主要倡议者相继离开中枢，改革推进中止”。其背景是改革触及考课、荐举和官僚利益，在中枢与地方均遭遇阻力，后续影响为仁宗朝保留部分制度讨论，但全面整饬转为未决议题。政权、军队或制度的变化须经由道路、文书、资源和地方中介才会落实，城市、乡村、边地和流动人群不可能在同一时间承受相同结果。

本条使用《宋史》卷10〈仁宗本纪三〉；《宋史》卷314〈范仲淹传〉；《续资治通鉴长编》庆历五年条，并标示“人事变化可靠；不能把离任简单解释为单一党争胜负”。这些材料能够钉住年序、命令、战役或局部地点，却不能直接推出人口、税收、身份和社会动机的总量。应区分诏令与实施、条约与边境生活、军报与伤亡统计，也不将官府的族群或行政称谓写成固定的社会本质。

从区域史角度看，该事件是资源、信息与权力重新连接的节点。不同地方会因市场、交通、宗教组织、家庭策略和环境条件而产生不同反应；材料的沉默限制我们对未被记录者所能作出的判断。

\eventanchor{NS-1045-REGNAL-YEAR}{1045年·宋以宋仁宗在位第23年作为诏令、赋役与外交文书的纪年框架}账本条目记录宋在1045年的“宋以宋仁宗在位第23年作为诏令、赋役与外交文书的纪年框架”。其背景是新岁颁历、年号和纪年是中枢与地方行政沟通的共同前提，后续影响为为同年政令、战事和地方材料提供日期锚点，不能单独推知具体政策执行。政权、军队或制度的变化须经由道路、文书、资源和地方中介才会落实，城市、乡村、边地和流动人群不可能在同一时间承受相同结果。

本条使用《宋史》卷1—23〈本纪〉；《续资治通鉴长编》本年条；《宋会要辑稿》相关门类，并标示“纪年框架可靠；该记录仅确证政权纪年连续，年内具体事项须由本年条另证”。这些材料能够钉住年序、命令、战役或局部地点，却不能直接推出人口、税收、身份和社会动机的总量。应区分诏令与实施、条约与边境生活、军报与伤亡统计，也不将官府的族群或行政称谓写成固定的社会本质。

从区域史角度看，该事件是资源、信息与权力重新连接的节点。不同地方会因市场、交通、宗教组织、家庭策略和环境条件而产生不同反应；材料的沉默限制我们对未被记录者所能作出的判断。

\eventanchor{NS-1046-REGNAL-YEAR}{1046年·宋以宋仁宗在位第24年作为诏令、赋役与外交文书的纪年框架}账本条目记录宋在1046年的“宋以宋仁宗在位第24年作为诏令、赋役与外交文书的纪年框架”。其背景是新岁颁历、年号和纪年是中枢与地方行政沟通的共同前提，后续影响为为同年政令、战事和地方材料提供日期锚点，不能单独推知具体政策执行。政权、军队或制度的变化须经由道路、文书、资源和地方中介才会落实，城市、乡村、边地和流动人群不可能在同一时间承受相同结果。

本条使用《宋史》卷1—23〈本纪〉；《续资治通鉴长编》本年条；《宋会要辑稿》相关门类，并标示“纪年框架可靠；该记录仅确证政权纪年连续，年内具体事项须由本年条另证”。这些材料能够钉住年序、命令、战役或局部地点，却不能直接推出人口、税收、身份和社会动机的总量。应区分诏令与实施、条约与边境生活、军报与伤亡统计，也不将官府的族群或行政称谓写成固定的社会本质。

从区域史角度看，该事件是资源、信息与权力重新连接的节点。不同地方会因市场、交通、宗教组织、家庭策略和环境条件而产生不同反应；材料的沉默限制我们对未被记录者所能作出的判断。

\eventanchor{NS-1047-REGNAL-YEAR}{1047年·宋以宋仁宗在位第25年作为诏令、赋役与外交文书的纪年框架}账本条目记录宋在1047年的“宋以宋仁宗在位第25年作为诏令、赋役与外交文书的纪年框架”。其背景是新岁颁历、年号和纪年是中枢与地方行政沟通的共同前提，后续影响为为同年政令、战事和地方材料提供日期锚点，不能单独推知具体政策执行。政权、军队或制度的变化须经由道路、文书、资源和地方中介才会落实，城市、乡村、边地和流动人群不可能在同一时间承受相同结果。

本条使用《宋史》卷1—23〈本纪〉；《续资治通鉴长编》本年条；《宋会要辑稿》相关门类，并标示“纪年框架可靠；该记录仅确证政权纪年连续，年内具体事项须由本年条另证”。这些材料能够钉住年序、命令、战役或局部地点，却不能直接推出人口、税收、身份和社会动机的总量。应区分诏令与实施、条约与边境生活、军报与伤亡统计，也不将官府的族群或行政称谓写成固定的社会本质。

从区域史角度看，该事件是资源、信息与权力重新连接的节点。不同地方会因市场、交通、宗教组织、家庭策略和环境条件而产生不同反应；材料的沉默限制我们对未被记录者所能作出的判断。

\eventanchor{NS-1047-WANG-ZE}{1047—1048年·贝州王则起事，宋军围城后平定}账本条目记录宋与河北起事者在1047—1048年的“贝州王则起事，宋军围城后平定”。其背景是河北边防、灾荒与基层控制压力为地方动员提供条件，后续影响为宋廷以军事平乱并检讨河北治理，个案不能外推为全国性社会运动。政权、军队或制度的变化须经由道路、文书、资源和地方中介才会落实，城市、乡村、边地和流动人群不可能在同一时间承受相同结果。

本条使用《宋史》卷10〈仁宗本纪三〉；《宋史》卷465〈王则传〉；《续资治通鉴长编》庆历七年至八年条，并标示“起事和平定可靠；宗教组织、参与者构成和死亡数字材料有限”。这些材料能够钉住年序、命令、战役或局部地点，却不能直接推出人口、税收、身份和社会动机的总量。应区分诏令与实施、条约与边境生活、军报与伤亡统计，也不将官府的族群或行政称谓写成固定的社会本质。

从区域史角度看，该事件是资源、信息与权力重新连接的节点。不同地方会因市场、交通、宗教组织、家庭策略和环境条件而产生不同反应；材料的沉默限制我们对未被记录者所能作出的判断。

\eventanchor{NS-1048-EVENT-001}{1048年·李元昊在内政冲突中遇害，西夏转入幼主与后族主政阶段}账本条目记录西夏与宋在1048年的“李元昊在内政冲突中遇害，西夏转入幼主与后族主政阶段”。其背景是前期边防、财政和统治联盟的压力构成直接背景，具体条件须按本条材料分辨，后续影响为改变后续资源配置与政治选择；实际执行范围和社会影响仍须分项核验。政权、军队或制度的变化须经由道路、文书、资源和地方中介才会落实，城市、乡村、边地和流动人群不可能在同一时间承受相同结果。

本条使用《宋史》卷1—23〈本纪〉；《续资治通鉴长编》本年条；《宋会要辑稿》相关门类，并标示“编年主干可靠；细节、规模与动机须与对方材料、地方文书或考古材料互校”。这些材料能够钉住年序、命令、战役或局部地点，却不能直接推出人口、税收、身份和社会动机的总量。应区分诏令与实施、条约与边境生活、军报与伤亡统计，也不将官府的族群或行政称谓写成固定的社会本质。

从区域史角度看，该事件是资源、信息与权力重新连接的节点。不同地方会因市场、交通、宗教组织、家庭策略和环境条件而产生不同反应；材料的沉默限制我们对未被记录者所能作出的判断。

\eventanchor{NS-1048-REGNAL-YEAR}{1048年·宋以宋仁宗在位第26年作为诏令、赋役与外交文书的纪年框架}账本条目记录宋在1048年的“宋以宋仁宗在位第26年作为诏令、赋役与外交文书的纪年框架”。其背景是新岁颁历、年号和纪年是中枢与地方行政沟通的共同前提，后续影响为为同年政令、战事和地方材料提供日期锚点，不能单独推知具体政策执行。政权、军队或制度的变化须经由道路、文书、资源和地方中介才会落实，城市、乡村、边地和流动人群不可能在同一时间承受相同结果。

本条使用《宋史》卷1—23〈本纪〉；《续资治通鉴长编》本年条；《宋会要辑稿》相关门类，并标示“纪年框架可靠；该记录仅确证政权纪年连续，年内具体事项须由本年条另证”。这些材料能够钉住年序、命令、战役或局部地点，却不能直接推出人口、税收、身份和社会动机的总量。应区分诏令与实施、条约与边境生活、军报与伤亡统计，也不将官府的族群或行政称谓写成固定的社会本质。

从区域史角度看，该事件是资源、信息与权力重新连接的节点。不同地方会因市场、交通、宗教组织、家庭策略和环境条件而产生不同反应；材料的沉默限制我们对未被记录者所能作出的判断。

\eventanchor{NS-1048-YELLOW-RIVER-SHANGHU}{1048年·黄河在商胡决口并改道北流，冲击河北水利与聚落}账本条目记录宋在1048年的“黄河在商胡决口并改道北流，冲击河北水利与聚落”。其背景是黄河堤防在高含沙量与长期维护压力下失守，后续影响为河患迫使宋廷持续投入河工和赈济，并改变河北交通、耕地和边防条件。政权、军队或制度的变化须经由道路、文书、资源和地方中介才会落实，城市、乡村、边地和流动人群不可能在同一时间承受相同结果。

本条使用《宋史》卷10〈仁宗本纪三〉；《宋会要辑稿·河渠》庆历八年河决条；黄河古河道沉积与堤防遗址研究，并标示“决口和改道可靠；河道持续时间、受灾人口和政策成效须以地层与地方材料检验”。这些材料能够钉住年序、命令、战役或局部地点，却不能直接推出人口、税收、身份和社会动机的总量。应区分诏令与实施、条约与边境生活、军报与伤亡统计，也不将官府的族群或行政称谓写成固定的社会本质。

从区域史角度看，该事件是资源、信息与权力重新连接的节点。不同地方会因市场、交通、宗教组织、家庭策略和环境条件而产生不同反应；材料的沉默限制我们对未被记录者所能作出的判断。

\eventanchor{NS-1049-REGNAL-YEAR}{1049年·宋以宋仁宗在位第27年作为诏令、赋役与外交文书的纪年框架}账本条目记录宋在1049年的“宋以宋仁宗在位第27年作为诏令、赋役与外交文书的纪年框架”。其背景是新岁颁历、年号和纪年是中枢与地方行政沟通的共同前提，后续影响为为同年政令、战事和地方材料提供日期锚点，不能单独推知具体政策执行。政权、军队或制度的变化须经由道路、文书、资源和地方中介才会落实，城市、乡村、边地和流动人群不可能在同一时间承受相同结果。

本条使用《宋史》卷1—23〈本纪〉；《续资治通鉴长编》本年条；《宋会要辑稿》相关门类，并标示“纪年框架可靠；该记录仅确证政权纪年连续，年内具体事项须由本年条另证”。这些材料能够钉住年序、命令、战役或局部地点，却不能直接推出人口、税收、身份和社会动机的总量。应区分诏令与实施、条约与边境生活、军报与伤亡统计，也不将官府的族群或行政称谓写成固定的社会本质。

从区域史角度看，该事件是资源、信息与权力重新连接的节点。不同地方会因市场、交通、宗教组织、家庭策略和环境条件而产生不同反应；材料的沉默限制我们对未被记录者所能作出的判断。

\eventanchor{NS-1050-REGNAL-YEAR}{1050年·宋以宋仁宗在位第28年作为诏令、赋役与外交文书的纪年框架}账本条目记录宋在1050年的“宋以宋仁宗在位第28年作为诏令、赋役与外交文书的纪年框架”。其背景是新岁颁历、年号和纪年是中枢与地方行政沟通的共同前提，后续影响为为同年政令、战事和地方材料提供日期锚点，不能单独推知具体政策执行。政权、军队或制度的变化须经由道路、文书、资源和地方中介才会落实，城市、乡村、边地和流动人群不可能在同一时间承受相同结果。

本条使用《宋史》卷1—23〈本纪〉；《续资治通鉴长编》本年条；《宋会要辑稿》相关门类，并标示“纪年框架可靠；该记录仅确证政权纪年连续，年内具体事项须由本年条另证”。这些材料能够钉住年序、命令、战役或局部地点，却不能直接推出人口、税收、身份和社会动机的总量。应区分诏令与实施、条约与边境生活、军报与伤亡统计，也不将官府的族群或行政称谓写成固定的社会本质。

从区域史角度看，该事件是资源、信息与权力重新连接的节点。不同地方会因市场、交通、宗教组织、家庭策略和环境条件而产生不同反应；材料的沉默限制我们对未被记录者所能作出的判断。

\eventanchor{NS-1051-REGNAL-YEAR}{1051年·宋以宋仁宗在位第29年作为诏令、赋役与外交文书的纪年框架}账本条目记录宋在1051年的“宋以宋仁宗在位第29年作为诏令、赋役与外交文书的纪年框架”。其背景是新岁颁历、年号和纪年是中枢与地方行政沟通的共同前提，后续影响为为同年政令、战事和地方材料提供日期锚点，不能单独推知具体政策执行。政权、军队或制度的变化须经由道路、文书、资源和地方中介才会落实，城市、乡村、边地和流动人群不可能在同一时间承受相同结果。

本条使用《宋史》卷1—23〈本纪〉；《续资治通鉴长编》本年条；《宋会要辑稿》相关门类，并标示“纪年框架可靠；该记录仅确证政权纪年连续，年内具体事项须由本年条另证”。这些材料能够钉住年序、命令、战役或局部地点，却不能直接推出人口、税收、身份和社会动机的总量。应区分诏令与实施、条约与边境生活、军报与伤亡统计，也不将官府的族群或行政称谓写成固定的社会本质。

从区域史角度看，该事件是资源、信息与权力重新连接的节点。不同地方会因市场、交通、宗教组织、家庭策略和环境条件而产生不同反应；材料的沉默限制我们对未被记录者所能作出的判断。

\eventanchor{NS-1052-NONG-ZHIGAO}{1052年·侬智高在岭南、交趾边境起兵并攻至邕州、广州附近}账本条目记录宋与广源州起事集团在1052年的“侬智高在岭南、交趾边境起兵并攻至邕州、广州附近”。其背景是边地首领在宋、交趾与地方部族网络间争取资源和承认，后续影响为宋廷调集南方军力，暴露岭南行政、道路和跨境治理的局限。政权、军队或制度的变化须经由道路、文书、资源和地方中介才会落实，城市、乡村、边地和流动人群不可能在同一时间承受相同结果。

本条使用《宋史》卷11〈仁宗本纪四〉；《宋史》卷495〈侬智高传〉；《续资治通鉴长编》皇祐四年条，并标示“起事路线大体可靠；族群身份和跨境归属不可按宋方标签简化”。这些材料能够钉住年序、命令、战役或局部地点，却不能直接推出人口、税收、身份和社会动机的总量。应区分诏令与实施、条约与边境生活、军报与伤亡统计，也不将官府的族群或行政称谓写成固定的社会本质。

从区域史角度看，该事件是资源、信息与权力重新连接的节点。不同地方会因市场、交通、宗教组织、家庭策略和环境条件而产生不同反应；材料的沉默限制我们对未被记录者所能作出的判断。

\eventanchor{NS-1052-REGNAL-YEAR}{1052年·宋以宋仁宗在位第30年作为诏令、赋役与外交文书的纪年框架}账本条目记录宋在1052年的“宋以宋仁宗在位第30年作为诏令、赋役与外交文书的纪年框架”。其背景是新岁颁历、年号和纪年是中枢与地方行政沟通的共同前提，后续影响为为同年政令、战事和地方材料提供日期锚点，不能单独推知具体政策执行。政权、军队或制度的变化须经由道路、文书、资源和地方中介才会落实，城市、乡村、边地和流动人群不可能在同一时间承受相同结果。

本条使用《宋史》卷1—23〈本纪〉；《续资治通鉴长编》本年条；《宋会要辑稿》相关门类，并标示“纪年框架可靠；该记录仅确证政权纪年连续，年内具体事项须由本年条另证”。这些材料能够钉住年序、命令、战役或局部地点，却不能直接推出人口、税收、身份和社会动机的总量。应区分诏令与实施、条约与边境生活、军报与伤亡统计，也不将官府的族群或行政称谓写成固定的社会本质。

从区域史角度看，该事件是资源、信息与权力重新连接的节点。不同地方会因市场、交通、宗教组织、家庭策略和环境条件而产生不同反应；材料的沉默限制我们对未被记录者所能作出的判断。

\eventanchor{NS-1053-DI-QING-GUANGNAN}{1053年·狄青率军平定侬智高起事}账本条目记录宋与广源州起事集团在1053年的“狄青率军平定侬智高起事”。其背景是侬智高控制区域交通并吸纳不满者，宋廷需跨区调兵，后续影响为宋恢复主要城镇控制，但边境社会的多重政治联系并未因此终结。政权、军队或制度的变化须经由道路、文书、资源和地方中介才会落实，城市、乡村、边地和流动人群不可能在同一时间承受相同结果。

本条使用《宋史》卷11〈仁宗本纪四〉；《宋史》卷290〈狄青传〉；《续资治通鉴长编》皇祐五年条，并标示“平定结果可靠；狄青个人战功与地方清剿范围多经传记放大”。这些材料能够钉住年序、命令、战役或局部地点，却不能直接推出人口、税收、身份和社会动机的总量。应区分诏令与实施、条约与边境生活、军报与伤亡统计，也不将官府的族群或行政称谓写成固定的社会本质。

从区域史角度看，该事件是资源、信息与权力重新连接的节点。不同地方会因市场、交通、宗教组织、家庭策略和环境条件而产生不同反应；材料的沉默限制我们对未被记录者所能作出的判断。

\eventanchor{NS-1053-REGNAL-YEAR}{1053年·宋以宋仁宗在位第31年作为诏令、赋役与外交文书的纪年框架}账本条目记录宋在1053年的“宋以宋仁宗在位第31年作为诏令、赋役与外交文书的纪年框架”。其背景是新岁颁历、年号和纪年是中枢与地方行政沟通的共同前提，后续影响为为同年政令、战事和地方材料提供日期锚点，不能单独推知具体政策执行。政权、军队或制度的变化须经由道路、文书、资源和地方中介才会落实，城市、乡村、边地和流动人群不可能在同一时间承受相同结果。

本条使用《宋史》卷1—23〈本纪〉；《续资治通鉴长编》本年条；《宋会要辑稿》相关门类，并标示“纪年框架可靠；该记录仅确证政权纪年连续，年内具体事项须由本年条另证”。这些材料能够钉住年序、命令、战役或局部地点，却不能直接推出人口、税收、身份和社会动机的总量。应区分诏令与实施、条约与边境生活、军报与伤亡统计，也不将官府的族群或行政称谓写成固定的社会本质。

从区域史角度看，该事件是资源、信息与权力重新连接的节点。不同地方会因市场、交通、宗教组织、家庭策略和环境条件而产生不同反应；材料的沉默限制我们对未被记录者所能作出的判断。

\eventanchor{NS-1054-REGNAL-YEAR}{1054年·宋以宋仁宗在位第32年作为诏令、赋役与外交文书的纪年框架}账本条目记录宋在1054年的“宋以宋仁宗在位第32年作为诏令、赋役与外交文书的纪年框架”。其背景是新岁颁历、年号和纪年是中枢与地方行政沟通的共同前提，后续影响为为同年政令、战事和地方材料提供日期锚点，不能单独推知具体政策执行。政权、军队或制度的变化须经由道路、文书、资源和地方中介才会落实，城市、乡村、边地和流动人群不可能在同一时间承受相同结果。

本条使用《宋史》卷1—23〈本纪〉；《续资治通鉴长编》本年条；《宋会要辑稿》相关门类，并标示“纪年框架可靠；该记录仅确证政权纪年连续，年内具体事项须由本年条另证”。这些材料能够钉住年序、命令、战役或局部地点，却不能直接推出人口、税收、身份和社会动机的总量。应区分诏令与实施、条约与边境生活、军报与伤亡统计，也不将官府的族群或行政称谓写成固定的社会本质。

从区域史角度看，该事件是资源、信息与权力重新连接的节点。不同地方会因市场、交通、宗教组织、家庭策略和环境条件而产生不同反应；材料的沉默限制我们对未被记录者所能作出的判断。

\eventanchor{NS-1055-REGNAL-YEAR}{1055年·宋以宋仁宗在位第33年作为诏令、赋役与外交文书的纪年框架}账本条目记录宋在1055年的“宋以宋仁宗在位第33年作为诏令、赋役与外交文书的纪年框架”。其背景是新岁颁历、年号和纪年是中枢与地方行政沟通的共同前提，后续影响为为同年政令、战事和地方材料提供日期锚点，不能单独推知具体政策执行。政权、军队或制度的变化须经由道路、文书、资源和地方中介才会落实，城市、乡村、边地和流动人群不可能在同一时间承受相同结果。

本条使用《宋史》卷1—23〈本纪〉；《续资治通鉴长编》本年条；《宋会要辑稿》相关门类，并标示“纪年框架可靠；该记录仅确证政权纪年连续，年内具体事项须由本年条另证”。这些材料能够钉住年序、命令、战役或局部地点，却不能直接推出人口、税收、身份和社会动机的总量。应区分诏令与实施、条约与边境生活、军报与伤亡统计，也不将官府的族群或行政称谓写成固定的社会本质。

从区域史角度看，该事件是资源、信息与权力重新连接的节点。不同地方会因市场、交通、宗教组织、家庭策略和环境条件而产生不同反应；材料的沉默限制我们对未被记录者所能作出的判断。

\eventanchor{NS-1056-REGNAL-YEAR}{1056年·宋以宋仁宗在位第34年作为诏令、赋役与外交文书的纪年框架}账本条目记录宋在1056年的“宋以宋仁宗在位第34年作为诏令、赋役与外交文书的纪年框架”。其背景是新岁颁历、年号和纪年是中枢与地方行政沟通的共同前提，后续影响为为同年政令、战事和地方材料提供日期锚点，不能单独推知具体政策执行。政权、军队或制度的变化须经由道路、文书、资源和地方中介才会落实，城市、乡村、边地和流动人群不可能在同一时间承受相同结果。

本条使用《宋史》卷1—23〈本纪〉；《续资治通鉴长编》本年条；《宋会要辑稿》相关门类，并标示“纪年框架可靠；该记录仅确证政权纪年连续，年内具体事项须由本年条另证”。这些材料能够钉住年序、命令、战役或局部地点，却不能直接推出人口、税收、身份和社会动机的总量。应区分诏令与实施、条约与边境生活、军报与伤亡统计，也不将官府的族群或行政称谓写成固定的社会本质。

从区域史角度看，该事件是资源、信息与权力重新连接的节点。不同地方会因市场、交通、宗教组织、家庭策略和环境条件而产生不同反应；材料的沉默限制我们对未被记录者所能作出的判断。

\eventanchor{NS-1057-EVENT-002}{1057年·嘉祐科举取士，苏轼、苏辙等士人进入中央政治视野}账本条目记录宋在1057年的“嘉祐科举取士，苏轼、苏辙等士人进入中央政治视野”。其背景是前期边防、财政和统治联盟的压力构成直接背景，具体条件须按本条材料分辨，后续影响为改变后续资源配置与政治选择；实际执行范围和社会影响仍须分项核验。政权、军队或制度的变化须经由道路、文书、资源和地方中介才会落实，城市、乡村、边地和流动人群不可能在同一时间承受相同结果。

本条使用《宋史》卷1—23〈本纪〉；《续资治通鉴长编》本年条；《宋会要辑稿》相关门类，并标示“编年主干可靠；细节、规模与动机须与对方材料、地方文书或考古材料互校”。这些材料能够钉住年序、命令、战役或局部地点，却不能直接推出人口、税收、身份和社会动机的总量。应区分诏令与实施、条约与边境生活、军报与伤亡统计，也不将官府的族群或行政称谓写成固定的社会本质。

从区域史角度看，该事件是资源、信息与权力重新连接的节点。不同地方会因市场、交通、宗教组织、家庭策略和环境条件而产生不同反应；材料的沉默限制我们对未被记录者所能作出的判断。
